% GENERATED FILE. Edit canonical lessons, then run npm run generate:books.
% canonical-source-hash: fnv1a64:aa996e50e0443d31
% canonical-lessons: PT-C397-dispositivo, PT-C397-portatil, PT-C397-instalar, PT-C397-desinstalar, PT-C397-atualizar

\chapter{Descarregar uma aplicação}
\label{ch:pt-descarregar-uma-aplicacao}
\hlchaptermodality{397}

\begin{chapteropening}
\textbf{By the end of this chapter:} \emph{I can install, update and uninstall apps on my laptop and my other devices.}

Complete the last lesson of chapter 397: I can install, update and uninstall apps on my laptop and my other devices.
\end{chapteropening}

\section[o dispositivo]{o dispositivo — the device}

\label{lesson:PT-C397-dispositivo}



\begin{quote}
Before the new one: say the Portuguese for the password, then the Portuguese for the profile.
\end{quote}

\subsection*{You'll want to know: o dispositivo}
\textbf{o dispositivo} — \textquotedblleft{}the device\textquotedblright{}.

\begin{quote}
Posso usar a aplicação em três dispositivos.
\end{quote}

\begin{quote}
\emph{I can use the app on three devices.}
\end{quote}

\emph{Dispositivo móvel} — mobile device.

\begin{cousinweb}
\emph{dispositivo} is a learned formation on Latin \emph{dispositus} \textquotedblleft{}arranged\textquotedblright{}, the participle of \emph{dispōnere} \textquotedblleft{}to set in order\textquotedblright{}, on \emph{pōnere} \textquotedblleft{}to put\textquotedblright{} (Lewis \& Short).
\end{cousinweb}

\begin{grammarlens}[title={What you've built}]
One more word for telling, explaining and describing at B1.
\end{grammarlens}

\subsection*{Your turn}
\begin{itemize}
\raggedright
  \item \emph{Say it:} \emph{o dispositivo}
  \item \emph{Say it:} \emph{o dispositivo}, once more
  \item \emph{Say it:} a sentence of your own with \emph{o dispositivo}
  \item \emph{Recall:} say the Portuguese for the password, then the Portuguese for the profile, then say \emph{o dispositivo} again
\end{itemize}

\begin{tcolorbox}[breakable,colback=teal!4,colframe=teal!35!black,title={Before you move on}]
What does o dispositivo mean? (\textquotedblleft{}The device\textquotedblright{}.) Say it once more.
\end{tcolorbox}
\section[o portátil]{o portátil — the laptop; portable}

\label{lesson:PT-C397-portatil}



\begin{quote}
Before the new one: say the Portuguese for the profile, then the Portuguese for the device.
\end{quote}

\subsection*{You'll want to know: o portátil}
\textbf{o portátil} — \textquotedblleft{}the laptop; portable\textquotedblright{}.

\begin{quote}
Levo o portátil para o trabalho todos os dias.
\end{quote}

\begin{quote}
\emph{I take my laptop to work every day.}
\end{quote}

European Portuguese; in Brazil \emph{o notebook}.

\begin{cousinweb}
\emph{portátil} comes from Medieval Latin \emph{portātilis} (EtymDB), built on \emph{portāre} \textquotedblleft{}to carry\textquotedblright{} (Lewis \& Short).
\end{cousinweb}

\begin{grammarlens}[title={What you've built}]
One more word for telling, explaining and describing at B1.
\end{grammarlens}

\subsection*{Your turn}
\begin{itemize}
\raggedright
  \item \emph{Say it:} \emph{o portátil}
  \item \emph{Say it:} \emph{o portátil}, once more
  \item \emph{Say it:} a sentence of your own with \emph{o portátil}
  \item \emph{Recall:} say the Portuguese for the profile, then the Portuguese for the device, then say \emph{o portátil} again
\end{itemize}

\begin{tcolorbox}[breakable,colback=teal!4,colframe=teal!35!black,title={Before you move on}]
What does o portátil mean? (\textquotedblleft{}The laptop; portable\textquotedblright{}.) Say it once more.
\end{tcolorbox}
\section[instalar]{instalar — to install}

\label{lesson:PT-C397-instalar}



\begin{quote}
Before the new one: say the Portuguese for the device, then the Portuguese for the laptop.
\end{quote}

\subsection*{You'll want to know: instalar}
\textbf{instalar} — \textquotedblleft{}to install\textquotedblright{}.

\begin{quote}
Instalei um jogo novo no portátil.
\end{quote}

\begin{quote}
\emph{I installed a new game on the laptop.}
\end{quote}

Also to settle in: \emph{instalar-se numa casa nova}.

\begin{cousinweb}
\emph{instalar} was borrowed from French \emph{installer} (EtymDB), which is usually traced to Medieval Latin \emph{installāre} \textquotedblleft{}to put into a stall (a seat in church)\textquotedblright{}, built on a Germanic word for \textquotedblleft{}standing place\textquotedblright{} (EtymDB lists its cousins).
\end{cousinweb}

\begin{grammarlens}[title={What you've built}]
One more word for telling, explaining and describing at B1.
\end{grammarlens}

\subsection*{Your turn}
\begin{itemize}
\raggedright
  \item \emph{Say it:} \emph{instalar}
  \item \emph{Say it:} \emph{instalar}, once more
  \item \emph{Say it:} a sentence of your own with \emph{instalar}
  \item \emph{Recall:} say the Portuguese for the device, then the Portuguese for the laptop, then say \emph{instalar} again
\end{itemize}

\begin{tcolorbox}[breakable,colback=teal!4,colframe=teal!35!black,title={Before you move on}]
What does instalar mean? (\textquotedblleft{}To install\textquotedblright{}.) Say it once more.
\end{tcolorbox}
\section[desinstalar]{desinstalar — to uninstall}

\label{lesson:PT-C397-desinstalar}



\begin{quote}
Before the new one: say the Portuguese for the laptop, then the Portuguese for to install.
\end{quote}

\subsection*{You'll want to know: desinstalar}
\textbf{desinstalar} — \textquotedblleft{}to uninstall\textquotedblright{}.

\begin{quote}
Desinstalei a aplicação porque não funcionava.
\end{quote}

\begin{quote}
\emph{I uninstalled the app because it did not work.}
\end{quote}

The opposite of \emph{instalar}.

\begin{cousinweb}
\emph{desinstalar} is built from \emph{des-} and \emph{instalar} \textquotedblleft{}to install\textquotedblright{}, which was borrowed from French \emph{installer} (EtymDB).
\end{cousinweb}

\begin{grammarlens}[title={What you've built}]
One more word for telling, explaining and describing at B1.
\end{grammarlens}

\subsection*{Your turn}
\begin{itemize}
\raggedright
  \item \emph{Say it:} \emph{desinstalar}
  \item \emph{Say it:} \emph{desinstalar}, once more
  \item \emph{Say it:} a sentence of your own with \emph{desinstalar}
  \item \emph{Recall:} say the Portuguese for the laptop, then the Portuguese for to install, then say \emph{desinstalar} again
\end{itemize}

\begin{tcolorbox}[breakable,colback=teal!4,colframe=teal!35!black,title={Before you move on}]
What does desinstalar mean? (\textquotedblleft{}To uninstall\textquotedblright{}.) Say it once more.
\end{tcolorbox}
\section[atualizar]{atualizar — to update}

\label{lesson:PT-C397-atualizar}



\begin{quote}
Before the new one: say the Portuguese for to install, then the Portuguese for to uninstall.
\end{quote}

\subsection*{You'll want to know: atualizar}
\textbf{atualizar} — \textquotedblleft{}to update\textquotedblright{}.

\begin{quote}
O telemóvel está a atualizar.
\end{quote}

\begin{quote}
\emph{The phone is updating.}
\end{quote}

AO90 \emph{atualizar}; older \emph{actualizar}. \emph{A atualização} — the update.

\begin{cousinweb}
\emph{atualizar} is built on \emph{atual} \textquotedblleft{}current\textquotedblright{}, from Late Latin \emph{āctuālis} \textquotedblleft{}active, practical\textquotedblright{}, on \emph{agere} \textquotedblleft{}to do\textquotedblright{} (Lewis \& Short).
\end{cousinweb}

\begin{grammarlens}[title={What you've built}]
One more word for telling, explaining and describing at B1.
\end{grammarlens}

\subsection*{Your turn}
\begin{itemize}
\raggedright
  \item \emph{Say it:} \emph{atualizar}
  \item \emph{Say it:} \emph{atualizar}, once more
  \item \emph{Say it:} a sentence of your own with \emph{atualizar}
  \item \emph{Recall:} say the Portuguese for to install, then the Portuguese for to uninstall, then say \emph{atualizar} again
\end{itemize}

\begin{tcolorbox}[breakable,colback=teal!4,colframe=teal!35!black,title={Before you move on}]
What does atualizar mean? (\textquotedblleft{}To update\textquotedblright{}.) Say it once more.
\end{tcolorbox}
